\documentclass[11pt,a4paper]{article}
\usepackage[utf8]{inputenc}
\usepackage[margin=1in]{geometry}
\usepackage{graphicx}
\usepackage{hyperref}
\usepackage{listings}
\usepackage{xcolor}
\usepackage{amsmath}
\usepackage{amssymb}

\lstset{
    language=C++,
    basicstyle=\ttfamily\small,
    keywordstyle=\color{blue},
    commentstyle=\color{gray},
    stringstyle=\color{red},
    breaklines=true,
    showstringspaces=false,
    frame=single,
    backgroundcolor=\color{lightgray!20}
}

\title{\textbf{Network Simulator Project} \\ \large{Weekly Progress Report}}
\author{Development Team}
\date{\today}

\begin{document}

\maketitle

\section{Executive Summary}

This week we extended the network simulator to include traffic generation, transport layer protocols (TCP/UDP), and network layer routing. These critical components form the foundation for realistic packet-level network simulation. All components have been implemented and tested, with comprehensive error documentation to guide future development.

\vspace{0.3cm}

\textbf{Components Built:}
\begin{itemize}
    \item Traffic Generation: \texttt{traffic\_generator.h / .cpp}
    \item Transport Protocols: \texttt{tcp.h / .cpp}, \texttt{udp.h / .cpp}
    \item Network Layer: \texttt{network\_layer.h / .cpp}
    \item Flow Management: \texttt{flow.h}
\end{itemize}

% ======================================================

\section{Phase 2: Traffic Generation, Transport \& Network Layer}

This section documents the critical issues encountered during implementation of traffic generation, transport layer protocols (TCP/UDP), and network layer routing, along with their solutions.

\subsection{1. Traffic Generation Layer}

\subsubsection{Error 1: Negative Traffic Rate}

\textbf{Problem:}
The system allowed Flow objects to be created with negative rates (e.g., \texttt{rate = -10}). The packet generation logic calculates the interval between successive packets as \texttt{interval = 1.0 / rate}. When rate is negative, this produces a negative interval, causing packets to be scheduled with past timestamps. This creates time paradoxes in the event scheduler and breaks the discrete event simulation model.

\textbf{Impact:}
The simulator's event queue becomes corrupted with negative timestamps. Packets may be scheduled before the current simulation time, causing the chronological ordering of events to violate causality. The simulation engine cannot process events out of order, resulting in either infinite loops or complete simulation failure.

\textbf{Fix:}
Validate the flow rate before scheduling:
\begin{lstlisting}
if (f.rate <= 0) {
    cerr << "ERROR: Flow rate must be positive\n";
    return;
}
\end{lstlisting}

\vspace{0.5cm}

\subsubsection{Error 2: Invalid Duration (Duration Smaller Than Offset)}

\textbf{Problem:}
Flows can be configured with a start time and duration such that the packet generation logic will immediately terminate the flow via the check \texttt{if (time > start\_time + duration) return;}. This causes the flow to generate very few packets despite a high rate setting.

\textbf{Impact:}
Traffic flows generate significantly fewer packets than intended, resulting in under-utilized network paths. The simulation produces misleading results about network load and packet delivery rates. Developers may not notice this error immediately since the code compiles without warnings and silently produces incorrect behavior.

\textbf{Fix:}
Add validation to ensure duration fields are reasonable:
\begin{lstlisting}
if (f.duration <= 0) {
    cerr << "ERROR: Duration must be positive\n";
    return;
}
if (f.start_time < 0) {
    cerr << "ERROR: Start time cannot be negative\n";
    return;
}
\end{lstlisting}

\vspace{0.5cm}

\subsubsection{Error 3: Zero Burst Rate}

\textbf{Problem:}
For BURST traffic patterns, the system allows configuration with \texttt{burst\_rate = 0}. The burst generation logic calculates inter-packet timing as \texttt{interval = 1.0 / burst\_rate}. When \texttt{burst\_rate = 0}, this causes division by zero, resulting in \texttt{interval = infinity} or \texttt{NaN}. These invalid floating-point values corrupt the event scheduler's timestamp field.

\textbf{Impact:}
Event timestamps become NaN or infinity, rendering the priority queue unable to sort events chronologically. The event scheduler fails to process events in the correct order, and the simulation either crashes with a floating-point exception or enters an infinite loop.

\textbf{Fix:}
Validate burst rate for BURST-type traffic:
\begin{lstlisting}
if (f.type == TrafficType::BURST && f.burst_rate <= 0) {
    cerr << "ERROR: Burst rate must be positive for BURST traffic\n";
    return;
}
\end{lstlisting}

\vspace{0.5cm}

\subsubsection{Error 4: Invalid Source Node}

\textbf{Problem:}
Flows are created with source or destination node IDs that don't exist in the network topology. The traffic generator calls \texttt{get\_node(source\_id)} to retrieve the source node object, which either crashes with a segmentation fault or returns a null/invalid pointer depending on the container implementation.

\textbf{Impact:}
The simulation crashes with undefined behavior. If the implementation uses a map, accessing a non-existent key may auto-create an invalid entry. Packets originating from invalid nodes corrupt the routing state and metrics collection.

\textbf{Fix:}
Validate node IDs before adding flows:
\begin{lstlisting}
if (f.source < 0 || f.source >= num_nodes) {
    cerr << "ERROR: Invalid source " << f.source << "\n";
    return;
}
if (f.destination < 0 || f.destination >= num_nodes) {
    cerr << "ERROR: Invalid destination " << f.destination << "\n";
    return;
}
\end{lstlisting}

\vspace{0.5cm}

\subsubsection{Error 5: Zero Duration Flow}

\textbf{Problem:}
A flow is created with \texttt{duration = 0}. The first packet is scheduled at \texttt{start\_time}. When the generation logic checks \texttt{if (time > start\_time + 0)}, the condition is immediately true. This causes at most one packet to be generated.

\textbf{Impact:}
Flows generate insufficient traffic, drastically reducing network load during simulation time periods. This produces unrealistic results for congestion, latency, and throughput measurements.

\textbf{Fix:}
Enforce positive duration:
\begin{lstlisting}
if (f.duration <= 0) {
    cerr << "ERROR: Duration must be > 0\n";
    return;
}
\end{lstlisting}

\vspace{0.5cm}

\subsection{2. Transport Layer (TCP/UDP)}

\subsubsection{Error 1: TCP State Uninitialized}

\textbf{Problem:}
TCP objects are created without explicit initialization of critical state variables. In C++, member variables have indeterminate values if not explicitly initialized in the constructor. When TCP::send() is called, it reads these uninitialized values to determine how many packets can be sent. The garbage values produce unpredictable congestion control behavior.

\textbf{Impact:}
Congestion control operates non-deterministically, making simulation results irreproducible and unreliable. Some simulation runs may pass packets freely while others drop all packets. Network load measurements become meaningless.

\textbf{Fix:}
Ensure all TCP state is properly initialized in the constructor:
\begin{lstlisting}
TCP::TCP() : cwnd(1), ssthresh(16), next_seq(1), 
             expected_ack(1), flow_id(0), 
             is_congestion_avoidance(false) { }
\end{lstlisting}

\vspace{0.5cm}

\subsubsection{Error 2: Receiving Expired Packets}

\textbf{Problem:}
UDP receives a packet with \texttt{TTL = 0}. The receive handler calls \texttt{record\_hop(node\_id)}, which decrements TTL to \texttt{-1}. The protocol layer does not check TTL validity before processing. A negative TTL violates the protocol specification.

\textbf{Impact:}
Expired and invalid packets are delivered to the application layer, corrupting metrics collection (hop counts become negative), violating end-to-end delivery semantics, and potentially causing memory errors.

\textbf{Fix:}
Check TTL validity before processing:
\begin{lstlisting}
void UDP::on_receive(const Packet& pkt) {
    if (pkt.ttl <= 0) {
        cout << "[UDP] Packet expired, dropping\n";
        return;
    }
}
\end{lstlisting}

\vspace{0.5cm}

\subsubsection{Error 3: TCP ACK Wraparound}

\textbf{Problem:}
TCP sequence numbers are 32-bit unsigned integers. When a flow runs for extended time, sequence numbers wrap around at $2^{32}$. If an ACK with maximum value arrives, the comparison may incorrectly allow acknowledgment of packets far beyond the current send window, corrupting the unacked packets map.

\textbf{Impact:}
The TCP state machine becomes inconsistent. Packets marked as acknowledged may not have been sent. Congestion window calculations become incorrect.

\textbf{Fix:}
Validate ACK sequence numbers:
\begin{lstlisting}
void TCP::on_ack(uint32_t ack_num) {
    if (ack_num > next_seq) {
        cerr << "ACK out of range: " << ack_num << "\n";
        return;
    }
}
\end{lstlisting}

\vspace{0.5cm}

\subsubsection{Error 4: Zero-Size Packets}

\textbf{Problem:}
Flow configuration allows \texttt{packet\_size = 0}. When packets with zero size traverse the link layer, transmission delay is calculated as \texttt{delay = (size * 8) / bandwidth}. With size = 0, transmission delay becomes 0 regardless of bandwidth.

\textbf{Impact:}
Network simulation produces incorrect timing because transmission time is not modeled. Link utilization calculations fail. Performance metrics become invalid.

\textbf{Fix:}
Validate packet size in flow configuration:
\begin{lstlisting}
if (f.packet_size <= 0) {
    cerr << "ERROR: Packet size must be > 0\n";
    return;
}
\end{lstlisting}

\vspace{0.5cm}

\subsection{3. Network Layer Routing}

\subsubsection{Error 1: Missing TTL Decrement}

\textbf{Problem:}
The NetworkLayer::receive() function processes incoming packets but forgets to decrement the TTL field before forwarding. Without TTL decrement, packets maintain the same TTL value at each hop. The TTL check never triggers, allowing packets to circulate indefinitely.

\textbf{Impact:}
Packets entering routing loops never terminate. Bandwidth becomes consumed by looping packets, congesting the network. Simulation time grows unbounded as events pile up from circulating packets.

\textbf{Fix:}
Always decrement TTL on packet reception:
\begin{lstlisting}
void NetworkLayer::receive(Packet p) {
    p.record_hop(node->id);
    p.ttl--;
    if (p.ttl <= 0) return;
    forward(p);
}
\end{lstlisting}

\vspace{0.5cm}

\subsubsection{Error 2: No Validation of Routing Entry}

\textbf{Problem:}
The routing table stores next-hop node IDs without validation. A developer can add a route pointing to a non-existent node. When forwarding a packet, the code calls \texttt{get\_node()}, which crashes or returns invalid memory.

\textbf{Impact:}
The network layer crashes when forwarding to routes pointing to non-existent nodes. Packets destined for unreachable destinations silently disappear. Network topology inconsistencies are not detected until runtime.

\textbf{Fix:}
Validate next hop before forwarding:
\begin{lstlisting}
if (next_hop == -1 || !nodes_map.count(next_hop)) {
    cerr << "Invalid route to " << next_hop << "\n";
    return;
}
\end{lstlisting}

\vspace{0.5cm}

\subsubsection{Error 3: Negative TTL Accepted}

\textbf{Problem:}
Packet constructor accepts TTL as a signed integer parameter, allowing negative values like \texttt{-5}. While negative TTL satisfy the check \texttt{if (ttl <= 0)}, allowing them violates the protocol specification.

\textbf{Impact:}
Invalid packets exist in the simulation, causing incorrect metrics (negative hop counts reported). Application layer code expecting non-negative TTL may crash. Test assertions will fail.

\textbf{Fix:}
Clamp TTL to valid range in packet constructor:
\begin{lstlisting}
Packet::Packet(..., int ttl_val) {
    ttl = max(0, ttl_val);
}
\end{lstlisting}

\vspace{0.5cm}

\subsubsection{Error 4: Self-Loop Packets (Source = Destination)}

\textbf{Problem:}
The network layer receives a packet with \texttt{source = destination}. The receive() function checks \texttt{if (p.destination == node->id)} to determine if the packet has reached its destination. This check returns true immediately, even if the packet was just generated locally with no network traversal.

\textbf{Impact:}
Self-destined flows produce unrealistic metrics: latency = 0, hop count = 0. This skews network statistics and bypasses the routing logic entirely.

\textbf{Fix:}
Optionally reject or warn on self-destined flows:
\begin{lstlisting}
if (f.source == f.destination) {
    cerr << "WARNING: Self-destined flow\n";
}
\end{lstlisting}

\vspace{0.5cm}

\subsubsection{Error 5: Cyclic Routes (Routing Loops)}

\textbf{Problem:}
Routing tables can be configured with cycles: Node 0 routes to B pointing to Node 1, and Node 1 routes to B pointing back to Node 0. The packet bounces between these nodes indefinitely, decrementing TTL by one each hop. It terminates only when TTL reaches zero after 64 hops.

\textbf{Impact:}
Routing loops waste network bandwidth with packets that will never reach their destination. Latency metrics become artificially high. Multiple looped flows cause severe network congestion.

\textbf{Fix:}
Implement loop detection or rely on TTL to break loops:
\begin{lstlisting}
// Recommended: Implement loop detection algorithm
// For now: TTL provides eventual termination.
\end{lstlisting}

\vspace{0.5cm}

\section{Test Coverage}

Three comprehensive test programs validate the error fixes:
\begin{itemize}
    \item \texttt{test/traffic\_generator\_test.cpp}: 5 traffic generation error scenarios
    \item \texttt{test/transport\_protocol\_test.cpp}: 4 transport protocol error scenarios
    \item \texttt{test/network\_layer\_test.cpp}: 5 network layer error scenarios
\end{itemize}

Each test file demonstrates the error condition and verifies that appropriate validation catches the issue before it causes simulation failure.

\section{Conclusion}

Phase 2 implementation added critical simulation layers: traffic injection, protocol handling, and routing. By documenting these 14 common pitfalls with detailed explanations and fixes, the simulator is now robust and extensible. The comprehensive test suite ensures future development maintains correctness and prevents regression.

\vfill
\hrule
\vspace{0.5cm}
\noindent
\textit{Report Updated: \today} \\
\textit{Project Status: Phase 2 Complete (Traffic, Transport, Network Layer)}

\end{document}
